\chapter{Optimiser, comparer et approcher}
\section*{Exercice 11.1 -- Boîte sans couvercle}
Le domaine physique est $0<x<3$. Le volume vaut $V(x)=x(6-2x)^2=4x(3-x)^2$. Sa dérivée est $12(x-1)(x-3)$, positive avant $1$ et négative entre $1$ et $3$. Le maximum est $V(1)=16$ unités cubes. Les valeurs limites aux extrémités sont nulles.

\section*{Exercice 11.2 -- Comparer deux expressions}
Algébriquement, $x+1/x-2=(x-1)^2/x\ge0$. Par dérivation, la dérivée vaut $1-1/x^2$, négative sur $]0,1[$ puis positive sur $]1,+\infty[$ ; le minimum est $2$ en $1$. L'hypothèse $x>0$ est indispensable au signe de la fraction.

\section*{Exercice 11.3 -- Erreur d'une tangente}
La tangente donne $1+3(0{,}02)=1{,}06$. Sur $[1,1{,}02]$, $f''(x)=6x\le6{,}12$. Le reste de Taylor est au plus $6{,}12(0{,}02)^2/2=0{,}001224$. La valeur exacte $1{,}061208$ confirme l'erreur $0{,}001208$, sans remplacer la preuve de la borne.

\section*{Exercice 11.4 -- Newton pour une racine}
$u_1=(2+2/2)/2=3/2$, puis $u_2=(3/2+4/3)/2=17/12$. Comme les itérés sont au-dessus de $\sqrt2$, $2/u_2=24/17\le\sqrt2\le17/12$. La largeur vaut $17/12-24/17=1/204$.

\section*{Exercice 11.5 -- Échec de Newton}
$f(0)=2$, $f'(0)=-2$ donnent $u_1=1$. Puis $f(1)=1$, $f'(1)=1$ donnent $u_2=0$. L'itération alterne entre $0$ et $1$ et ne converge pas. Aucun dénominateur n'est nul : vérifier seulement ce point ne suffit pas à garantir la convergence.

\section*{Exercice 11.6 -- Maximum sur une grille}
La fonction atteint un maximum $f(x_*)$ sur le segment. Un point de grille $g$ est à distance au plus $0{,}01$ de $x_*$. Les accroissements finis donnent $f(x_*)-f(g)\le5\times0{,}01=0{,}05$. Le maximum sur la grille est au moins $f(g)$ et au plus le maximum vrai : l'écart est donc entre $0$ et $0{,}05$.
